\honest{Superstimuli and the Collapse of Western Civilization}{Eliezer Yudkowsky}{2007}

At least three people have died playing online games for days without rest. \nb{The three linked reports match real cases: a Chinese World of Warcraft player in 2005, a South Korean after about fifty hours of StarCraft in 2005, and a Chinese man after a seven-day session in 2007.} If people have the right to play games, the market will supply the most engaging games that can be sold, even to the point of removing the most engaged consumers from the gene pool. How does a product become so involving that after 57 hours a person would rather use it for one more hour than eat or sleep?

A candy bar is a ``superstimulus'': it matches taste buds that evolved among hunter-gatherers more strongly than anything they ever met. Tastiness, which once tracked healthy food, has been reverse-engineered and matched artificially; there is no equal market incentive to make the food healthy, since we cannot taste health. \nb{The idea is Niko Tinbergen's ``supernormal stimulus,'' from experiments in which birds preferred oversized dummy eggs to their own. The post does not mention Tinbergen.} The billboard model in Dove's ``Evolution'' video is another, a beauty no real woman can match. Real women are killing themselves to keep up, supermodels for example with cocaine. I give no source. A video game can be more engaging than reality; I guess at the tricks: challenges poised between easy and impossible, intermittent reinforcement, an ever-rising score, and social play in multiplayer games.

One might hope for a sweet spot where most players have fun and only a few become so addicted they lose their jobs and cannot pay; in 2007, playing 58 hours until you die is still the exception. But competition will not stop there: ``if you can make your game 5\% more addictive, you may be able to steal 50\% of your competitor's customers.'' The market will supply as much temptation as can be sold, beyond the point where stimuli become superstimuli (consider how advertised beauty has changed since the 1950s) and, as candy bars show, past the point of harm.

Why not just say no? Free-market economics assumes that, absent force and fraud, people can always refuse a harmful transaction; to the extent that is true, the free market has few downsides. An organism that regularly passes up food dies, but sometimes a usually beneficial act is harmful, and humans can see such cases by abstract thought, though we also imagine ones that do not exist, like ancestor spirits forbidding us good rabbits. Evolution struck a compromise, a limited ability to resist, since our ancestors with too much willpower ``probably starved themselves to sacrifice to the gods.'' And ``Resisting any temptation takes conscious expenditure of an exhaustible supply of mental energy.'' \nb{The link is to an economic model that assumes this. The psychology behind it, ego depletion, later failed two large preregistered replications, in 2016 and 2021.} The players who died must, in a sense, have used willpower to keep playing. Even people lucky in willpower pay a price to resist; it is just easier for them to pay. Willpower evolved against ancestral temptations, and may run out faster, ``it seems plausible,'' against superstimuli.

Is the public display of superstimuli a harm even to those who resist? Should we ban cookie ads? A problem does not prove that government can fix it, since regulators focus on failures spectacular enough to reach the newspapers, not on low-grade harm from addictive products; and government's failure does not prove that nothing is wrong.

Last comes ``a final argument from fictional evidence,'' Simon Funk's novel \textsc{After Life}, in which humans are replaced by artificial children more fun to raise than real ones. ``Perhaps'' birth rates in rich countries are falling because the market supplies ever more tempting alternatives to children. ``Where are the advertising billboards that say `BREED'?'' Who will pay image consultants to make arguing with sullen teenagers seem more alluring than a vacation in Tahiti? My last line quotes Funk: ``In the end,'' the human species ``was simply marketed out of existence.'' \nb{That sentence could not be found in the text of the novel that Funk offers for download. The nearest line there is a character's question, ``So, we simply marketed the humans out of existence?''}
